% GENERATED FILE. Edit canonical lessons, then run npm run generate:books.
% canonical-source-hash: fnv1a64:b18a319667448fa3
% canonical-lessons: MW-C276-basli, MW-C276-bajavno, MW-C276-sur, MW-C276-tali, MW-C276-gujno

\chapter{Flute and Clapping}
\label{ch:mw-flute-and-clapping}
\hlchaptermodality{276}

\begin{chapteropening}
\textbf{By the end of this chapter:} \emph{I can say which music I like to play.}

Complete the last lesson of chapter 276: I can say which music I like to play.
\end{chapteropening}

\section[\texorpdfstring{\mw{बांसळी} (bā̃sḷī)}{बांसळी (bā̃sḷī)}]{\mw{बांसळी} (bā̃sḷī) — a flute}

\label{lesson:MW-C276-basli}



\begin{quote}
Before the new one: say the Marwadi for credit for a deed, then the Marwadi for an opinion.
\end{quote}

\subsection*{You'll want to know: \mw{बांसळी}}
\textbf{\mw{बांसळी}} — \emph{bā̃sḷī} — \textquotedblleft{}a flute\textquotedblright{}.

\begin{quote}
\mw{म्हारी} \mw{बांसळी} \mw{लांबी} \mw{है।}
\end{quote}

\begin{quote}
\emph{mhārī bā̃sḷī lā̃bī hai.}
\end{quote}

\begin{quote}
My flute is long.
\end{quote}

A feminine noun.

\begin{cousinweb}
\textbf{\mw{बांसळी}} (\emph{bā̃sḷī}) is inherited from Sanskrit \emph{vamśī}, \textquotedblleft{}a flute, a pipe\textquotedblright{}, itself from \emph{vamśa}, \textquotedblleft{}bamboo\textquotedblright{}; CDIAL gives no Marwari form in this entry, but lists Gujarati \emph{vā̃sḷī}, \textquotedblleft{}pipe, flute\textquotedblright{}, and Hindi \emph{bãsī}, \textquotedblleft{}flute\textquotedblright{} (CDIAL 11180).
\end{cousinweb}

\begin{grammarlens}[title={What you've built}]
One more word for this chapter.
\end{grammarlens}

\subsection*{Your turn}
\begin{itemize}
\raggedright
  \item \emph{Say it:} \emph{bā̃sḷī}
  \item \emph{Say it:} \emph{bā̃sḷī}, once more
  \item \emph{Say it:} \emph{mhārī bā̃sḷī lā̃bī hai.} — My flute is long.
  \item \emph{Recall:} say the Marwadi for credit for a deed, then the Marwadi for an opinion, then say \emph{bā̃sḷī} again
\end{itemize}

\begin{tcolorbox}[breakable,colback=teal!4,colframe=teal!35!black,title={Before you move on}]
What does \mw{बांसळी} mean? (\textquotedblleft{}A flute\textquotedblright{}.) Say it once more.
\end{tcolorbox}
\section[\texorpdfstring{\mw{बजावणो} (bajāvṇo)}{बजावणो (bajāvṇo)}]{\mw{बजावणो} (bajāvṇo) — to play an instrument}

\label{lesson:MW-C276-bajavno}



\begin{quote}
Before the new one: say the Marwadi for an opinion, then the Marwadi for a flute.
\end{quote}

\subsection*{You'll want to know: \mw{बजावणो}}
\textbf{\mw{बजावणो}} — \emph{bajāvṇo} — \textquotedblleft{}to play an instrument\textquotedblright{}.

\begin{quote}
\mw{भाई} \mw{बांसळी} \mw{बजावै} \mw{है।}
\end{quote}

\begin{quote}
\emph{bhāī bā̃sḷī bajāvai hai.}
\end{quote}

\begin{quote}
My brother plays the flute.
\end{quote}

Past (I …): \mw{म्हैं} \mw{बांसळी} \mw{बजाई}, \emph{mhaĩ bā̃sḷī bajāī}.

\begin{cousinweb}
\textbf{\mw{बजावणो}} (\emph{bajāvṇo}) is built on the verb CDIAL lists as Old Marwari \emph{bājaï}, \textquotedblleft{}to sound\textquotedblright{}, under Sanskrit \emph{vādyate}, \textquotedblleft{}is made to sound\textquotedblright{} (CDIAL 11513); as a causative it means \textquotedblleft{}to make sound\textquotedblright{}.
\end{cousinweb}

\begin{grammarlens}[title={What you've built}]
One more word for this chapter.
\end{grammarlens}

\subsection*{Your turn}
\begin{itemize}
\raggedright
  \item \emph{Say it:} \emph{bajāvṇo}
  \item \emph{Say it:} \emph{bajāvṇo}, once more
  \item \emph{Say it:} \emph{bhāī bā̃sḷī bajāvai hai.} — My brother plays the flute.
  \item \emph{Recall:} say the Marwadi for an opinion, then the Marwadi for a flute, then say \emph{bajāvṇo} again
\end{itemize}

\begin{tcolorbox}[breakable,colback=teal!4,colframe=teal!35!black,title={Before you move on}]
What does \mw{बजावणो} mean? (\textquotedblleft{}To play an instrument\textquotedblright{}.) Say it once more.
\end{tcolorbox}
\section[\texorpdfstring{\mw{सुर} (sur)}{सुर (sur)}]{\mw{सुर} (sur) — a tune}

\label{lesson:MW-C276-sur}



\begin{quote}
Before the new one: say the Marwadi for a flute, then the Marwadi for to play an instrument.
\end{quote}

\subsection*{You'll want to know: \mw{सुर}}
\textbf{\mw{सुर}} — \emph{sur} — \textquotedblleft{}a tune\textquotedblright{}.

\begin{quote}
\mw{यो} \mw{सुर} \mw{मीठो} \mw{है।}
\end{quote}

\begin{quote}
\emph{yo sur mīṭho hai.}
\end{quote}

\begin{quote}
This tune is sweet.
\end{quote}

A masculine noun.

\begin{cousinweb}
\textbf{\mw{सुर}} (\emph{sur}) is inherited from Sanskrit \emph{sura}, \textquotedblleft{}sound\textquotedblright{}, a variant of \emph{svara}; CDIAL gives no Marwari form in this entry, but lists Hindi \emph{sur}, \textquotedblleft{}sound, words, tune\textquotedblright{} (CDIAL 13498).
\end{cousinweb}

\begin{grammarlens}[title={What you've built}]
One more word for this chapter.
\end{grammarlens}

\subsection*{Your turn}
\begin{itemize}
\raggedright
  \item \emph{Say it:} \emph{sur}
  \item \emph{Say it:} \emph{sur}, once more
  \item \emph{Say it:} \emph{yo sur mīṭho hai.} — This tune is sweet.
  \item \emph{Recall:} say the Marwadi for a flute, then the Marwadi for to play an instrument, then say \emph{sur} again
\end{itemize}

\begin{tcolorbox}[breakable,colback=teal!4,colframe=teal!35!black,title={Before you move on}]
What does \mw{सुर} mean? (\textquotedblleft{}A tune\textquotedblright{}.) Say it once more.
\end{tcolorbox}
\section[\texorpdfstring{\mw{ताळी} (tāḷī)}{ताळी (tāḷī)}]{\mw{ताळी} (tāḷī) — a clap}

\label{lesson:MW-C276-tali}



\begin{quote}
Before the new one: say the Marwadi for to play an instrument, then the Marwadi for a tune.
\end{quote}

\subsection*{You'll want to know: \mw{ताळी}}
\textbf{\mw{ताळी}} — \emph{tāḷī} — \textquotedblleft{}a clap\textquotedblright{}.

\begin{quote}
\mw{सगळा} \mw{ताळी} \mw{बजावै} \mw{है।}
\end{quote}

\begin{quote}
\emph{sagḷā tāḷī bajāvai hai.}
\end{quote}

\begin{quote}
Everyone claps.
\end{quote}

A feminine noun.

\begin{cousinweb}
\textbf{\mw{ताळी}} (\emph{tāḷī}) is inherited from Sanskrit \emph{tāla}, \textquotedblleft{}a clap, rhythm\textquotedblright{}; CDIAL gives no Marwari form in this entry, but lists Gujarati \emph{tāḷī}, \textquotedblleft{}clapping\textquotedblright{} (CDIAL 5748). You clap with \textbf{\mw{बजावणो}}: \textbf{\mw{ताळी} \mw{बजावो}}.
\end{cousinweb}

\begin{grammarlens}[title={What you've built}]
One more word for this chapter.
\end{grammarlens}

\subsection*{Your turn}
\begin{itemize}
\raggedright
  \item \emph{Say it:} \emph{tāḷī}
  \item \emph{Say it:} \emph{tāḷī}, once more
  \item \emph{Say it:} \emph{sagḷā tāḷī bajāvai hai.} — Everyone claps.
  \item \emph{Recall:} say the Marwadi for to play an instrument, then the Marwadi for a tune, then say \emph{tāḷī} again
\end{itemize}

\begin{tcolorbox}[breakable,colback=teal!4,colframe=teal!35!black,title={Before you move on}]
What does \mw{ताळी} mean? (\textquotedblleft{}A clap\textquotedblright{}.) Say it once more.
\end{tcolorbox}
\section[\texorpdfstring{\mw{गूंजणो} (gū̃jṇo)}{गूंजणो (gū̃jṇo)}]{\mw{गूंजणो} (gū̃jṇo) — to hum; to ring out}

\label{lesson:MW-C276-gujno}



\begin{quote}
Before the new one: say the Marwadi for a tune, then the Marwadi for a clap.
\end{quote}

\subsection*{You'll want to know: \mw{गूंजणो}}
\textbf{\mw{गूंजणो}} — \emph{gū̃jṇo} — \textquotedblleft{}to hum; to ring out\textquotedblright{}.

\begin{quote}
\mw{मंदिर} \mw{मांय} \mw{गीत} \mw{गूंजै} \mw{है।}
\end{quote}

\begin{quote}
\emph{mandir mā̃y gīt gū̃jai hai.}
\end{quote}

\begin{quote}
A song rings out in the temple.
\end{quote}

Past (it …): \mw{वो} \mw{गीत} \mw{गूंज्यो}, \emph{vo gīt gū̃jyo}.

\begin{cousinweb}
\textbf{\mw{गूंजणो}} (\emph{gū̃jṇo}) is inherited from Sanskrit \emph{guñjati}, \textquotedblleft{}hums, buzzes\textquotedblright{}; CDIAL gives no Marwari form in this entry, but lists Hindi \emph{gū̃jnā}, \textquotedblleft{}to hum\textquotedblright{} (CDIAL 4175). A sound that hums on fills the room.
\end{cousinweb}

\begin{grammarlens}[title={What you've built}]
One more word for this chapter.
\end{grammarlens}

\subsection*{Your turn}
\begin{itemize}
\raggedright
  \item \emph{Say it:} \emph{gū̃jṇo}
  \item \emph{Say it:} \emph{gū̃jṇo}, once more
  \item \emph{Say it:} \emph{mandir mā̃y gīt gū̃jai hai.} — A song rings out in the temple.
  \item \emph{Recall:} say the Marwadi for a tune, then the Marwadi for a clap, then say \emph{gū̃jṇo} again
\end{itemize}

\begin{tcolorbox}[breakable,colback=teal!4,colframe=teal!35!black,title={Before you move on}]
What does \mw{गूंजणो} mean? (\textquotedblleft{}To hum; to ring out\textquotedblright{}.) Say it once more.
\end{tcolorbox}
